\documentclass[12pt]{article}
\usepackage[margin=1in]{geometry}
\usepackage{amsmath,booktabs,graphicx,makecell,threeparttable}

\begin{document}

\providecommand{\StructuralRobustnessPath}{appendix/structural-robustness}

\section{Prior sensitivity of the social multiplier}
\label{app:multiplier-prior-sensitivity}

We assess whether the estimated level and treatment ordering of the social
multiplier are driven by the maintained priors.  The exercise varies five
modules separately: the social-image weight, the direct distance cost, the
scale of idiosyncratic heterogeneity, the first-order visibility schedule, and
private utility (including the WTP link).  For each module we halve and double
the baseline scale, where applicable.  The inverse-gamma alternatives for the
heterogeneity scale hold its prior mean fixed, so that they vary concentration
rather than location.  We also report joint tight and joint diffuse stress
tests that move all five modules together.

Concretely, the baseline/tight/diffuse scales are $1/0.5/2$ for the
social-image weight, $0.25/0.125/0.5$ for the direct distance coefficient,
and $1/0.5/2$ for the visibility intercepts (with the corresponding contrast
scales always one half as large).  The inverse-gamma shape--scale priors for
idiosyncratic heterogeneity are $(3.3,1.1)$, $(8,3.348)$, and
$(2.5,0.717)$;
all have mean 0.478.  For private utility, the baseline/tight/diffuse scales
are $1/0.5/2$ for the intercept, $0.25/0.125/0.5$ for treatment effects, and
$2/1/4$ for the WTP location.  The positive WTP-to-utility link uses
$N^+(0,10^{-4})$, $N^+(0,5\times10^{-5})$, and, as a deliberately broad
tail stress test, $\operatorname{lognormal}(-10,4)$.

This is an ordinary-posterior prior sensitivity exercise, separate from the
preferred exponential community-weight analysis in
Section~\ref{app:cluster-inference}.  Every specification uses the same data,
likelihood, common social-image weight, and Gaussian type distribution as the
main structural model.  The primary estimand is the finite multiplier over
0--2.5 km, which summarizes the economically relevant function over one
prespecified distance span rather than making simultaneous claims at each
point of a dense grid.  In addition to arm-specific levels and 95 percent
equal-tailed credible intervals, we report the posterior probabilities that
the average No Signal multiplier exceeds the average Signal multiplier and
that the Calendar multiplier exceeds the Bracelet multiplier.

\IfFileExists{\StructuralRobustnessPath/tables/main-core-prior-grid.tex}{%
\begin{table}[htbp]
  \centering
  \caption{Prior sensitivity of the finite social multiplier}
  \label{tab:main-core-prior-grid}
  \begin{threeparttable}
    \centering
\resizebox{\ifdim\width>\linewidth\linewidth\else\width\fi}{!}{%
\begin{tabular}[t]{lcccccc}
\toprule
 & \multicolumn{4}{c}{Finite multiplier, 0--2.5 km} & \multicolumn{2}{c}{Directional posterior probability} \\
\cmidrule(l{3pt}r{3pt}){2-5}\cmidrule(l{3pt}r{3pt}){6-7}
Prior specification & Bracelet & Calendar & Ink & Control & $M_{NS}>M_S$ & $M_C>M_B$ \\
\midrule
Baseline & \makecell[c]{0.84\\(0.70, 0.97)} & \makecell[c]{1.23\\(1.07, 1.48)} & \makecell[c]{0.88\\(0.76, 0.99)} & \makecell[c]{1.47\\(1.24, 1.84)} & >0.999 & >0.999\\
\addlinespace
\multicolumn{7}{l}{\textit{Social-image weight}}\\
\addlinespace
Tight & \makecell[c]{0.86\\(0.72, 0.98)} & \makecell[c]{1.21\\(1.06, 1.42)} & \makecell[c]{0.89\\(0.79, 0.98)} & \makecell[c]{1.43\\(1.22, 1.74)} & >0.999 & >0.999\\
\addlinespace
Diffuse & \makecell[c]{0.84\\(0.70, 0.98)} & \makecell[c]{1.24\\(1.07, 1.50)} & \makecell[c]{0.88\\(0.75, 0.98)} & \makecell[c]{1.49\\(1.25, 1.84)} & >0.999 & >0.999\\
\addlinespace
\multicolumn{7}{l}{\textit{Distance cost}}\\
\addlinespace
Tight & \makecell[c]{0.85\\(0.70, 0.97)} & \makecell[c]{1.24\\(1.08, 1.49)} & \makecell[c]{0.88\\(0.76, 1.00)} & \makecell[c]{1.49\\(1.26, 1.89)} & >0.999 & >0.999\\
\addlinespace
Diffuse & \makecell[c]{0.85\\(0.70, 0.98)} & \makecell[c]{1.23\\(1.07, 1.49)} & \makecell[c]{0.88\\(0.75, 0.99)} & \makecell[c]{1.47\\(1.24, 1.82)} & >0.999 & >0.999\\
\addlinespace
\multicolumn{7}{l}{\textit{Idiosyncratic heterogeneity}}\\
\addlinespace
Tight & \makecell[c]{0.85\\(0.69, 0.98)} & \makecell[c]{1.23\\(1.08, 1.47)} & \makecell[c]{0.88\\(0.76, 1.00)} & \makecell[c]{1.47\\(1.24, 1.85)} & >0.999 & 0.999\\
\addlinespace
Diffuse & \makecell[c]{0.84\\(0.69, 0.98)} & \makecell[c]{1.24\\(1.07, 1.49)} & \makecell[c]{0.88\\(0.75, 0.99)} & \makecell[c]{1.48\\(1.25, 1.82)} & >0.999 & >0.999\\
\addlinespace
\multicolumn{7}{l}{\textit{Visibility schedule}}\\
\addlinespace
Tight & \makecell[c]{0.86\\(0.71, 0.98)} & \makecell[c]{1.24\\(1.07, 1.49)} & \makecell[c]{0.88\\(0.76, 0.98)} & \makecell[c]{1.40\\(1.19, 1.73)} & >0.999 & >0.999\\
\addlinespace
Diffuse & \makecell[c]{0.84\\(0.68, 0.97)} & \makecell[c]{1.23\\(1.06, 1.47)} & \makecell[c]{0.88\\(0.76, 0.99)} & \makecell[c]{1.51\\(1.27, 1.89)} & >0.999 & >0.999\\
\addlinespace
\multicolumn{7}{l}{\textit{Private utility}}\\
\addlinespace
Tight & \makecell[c]{0.86\\(0.74, 0.97)} & \makecell[c]{1.19\\(1.06, 1.38)} & \makecell[c]{0.90\\(0.79, 0.98)} & \makecell[c]{1.39\\(1.20, 1.69)} & >0.999 & >0.999\\
\addlinespace
Diffuse & \makecell[c]{0.84\\(0.68, 0.98)} & \makecell[c]{1.25\\(1.07, 1.51)} & \makecell[c]{0.87\\(0.74, 0.99)} & \makecell[c]{1.52\\(1.27, 1.89)} & >0.999 & >0.999\\
\addlinespace
\multicolumn{7}{l}{\textit{Joint stress test}}\\
\addlinespace
Tight & \makecell[c]{0.89\\(0.78, 0.97)} & \makecell[c]{1.18\\(1.07, 1.35)} & \makecell[c]{0.91\\(0.81, 0.99)} & \makecell[c]{1.30\\(1.14, 1.53)} & >0.999 & >0.999\\
\addlinespace
Diffuse & \makecell[c]{0.83\\(0.66, 0.97)} & \makecell[c]{1.25\\(1.08, 1.52)} & \makecell[c]{0.87\\(0.75, 0.99)} & \makecell[c]{1.55\\(1.29, 1.98)} & >0.999 & >0.999\\
\bottomrule
\end{tabular}%
}

    \begin{tablenotes}[flushleft]\footnotesize
      \item Notes: Entries in the four arm columns are posterior medians with
      95 percent equal-tailed credible intervals in parentheses.  The last two
      columns are directional posterior probabilities.  The finite multiplier
      is calculated over 0--2.5 km.  Tight and diffuse specifications change
      only the named prior module; the joint rows change all modules together.
    \end{tablenotes}
  \end{threeparttable}
\end{table}
}{}

The result is highly stable within this maintained likelihood.  Across all 13
specifications, posterior medians range from 0.83 to 0.89 for Bracelet, 0.87
to 0.91 for Ink, 1.18 to 1.25 for Calendar, and 1.30 to 1.55 for Control.  The
smallest posterior probability that Calendar exceeds Bracelet is 0.999, and
the posterior probability that the average No Signal multiplier exceeds the
average Signal multiplier is greater than 0.999 in every row.  Thus the level
and ordering of the finite multiplier are not artifacts of any one baseline
prior module.  This exercise is conditional on the baseline observation
model, however, and does not undo the separate sensitivity to asymmetric
reporting documented in Section~\ref{app:asymmetric-reporting-priors}.

\end{document}
